\begin{pastexam}[label={px:ncu-106-dsa-04}]{中央資工所 106 學年度 第 4 題}
（問答）

(10 points) Fill the following two blanks to complete an insertion sort algorithm (which sorts
elements of array $A[1{:}N]$).
\begin{minted}[fontsize=\small,breaklines]{c}
void insertionSort(int A[], int N)
{  int j, t;
   For (j=2; j<=N; j++){
        t = A[j];
        insert(t, A, j-1);}
}
void insert(int t, int A[], int i)
{
    A[0]=t;
    While (t<A[i]) {
        A[i+1] =__________ ;
        i--;}
    A[i+1] = __________;
}
\end{minted}
\end{pastexam}

\begin{pastsoln}
答案：第一格 \texttt{A[i]}，第二格 \texttt{t}。

\textbf{\texttt{insert} 在做什麼}。 呼叫 \texttt{insert(t, A, j-1)} 時，$A[1..j-1]$ 已經排好，
$t$ 是原來的 $A[j]$。 \texttt{insert} 要把 $t$ 插進 $A[1..i]$（$i = j - 1$）的正確位置：從右往左看，
比 $t$ 大的數都往右挪一格，遇到不比 $t$ 大的數就停，$t$ 放進它右邊空出來的那一格。 這就是資結書
第~1.3~節的插入排序（虛擬碼~1.3），只是把內層迴圈拆成一個函式。

\begin{itemize}
\item 迴圈裡：條件 $t < A[i]$ 成立表示 $A[i]$ 比 $t$ 大，要往右挪一格，所以
\texttt{A[i+1] = A[i]}，再 \texttt{i--} 看左邊下一個。 原來的 $A[j]$ 已經存在 $t$ 裡，第一次挪動
蓋掉 $A[j]$ 不會弄丟資料。
\item 迴圈結束時 $A[i] \le t$，而 $A[i+1]$ 已經挪到右邊（或本來就是 $t$ 原來的位置），是空出來的
那一格，所以 \texttt{A[i+1] = t}。
\end{itemize}

\textbf{$A[0] = t$ 的用處}。 這是資結書第~1.3.1~節提到的\textbf{哨兵}。 $t$ 比前面全部的數都小時，
$i$ 會一路減到 $0$，這時 $t < A[0] = t$ 不成立，迴圈一定停下來，所以條件不必再寫 $i \ge 1$。
停在 $i = 0$ 時，第二格把 $t$ 放進 $A[1]$，正確。

\textbf{追蹤}。 以 $A[1..6] = 5, 2, 4, 6, 1, 3$ 為例，每一列是 \texttt{insert} 做完之後：
\begin{center}
\begin{tabular}{cccl}
\toprule
$j$ & $t$ & 挪動次數，停下時的 $i$ & $A[1..6]$ \\
\midrule
2 & 2 & 1 次，$i = 0$（停在哨兵） & 2 5 4 6 1 3 \\
3 & 4 & 1 次，$i = 1$（$A[1] = 2 \le 4$） & 2 4 5 6 1 3 \\
4 & 6 & 0 次，$i = 3$（$A[3] = 5 \le 6$） & 2 4 5 6 1 3 \\
5 & 1 & 4 次，$i = 0$（停在哨兵） & 1 2 4 5 6 3 \\
6 & 3 & 3 次，$i = 2$（$A[2] = 2 \le 3$） & 1 2 3 4 5 6 \\
\bottomrule
\end{tabular}
\end{center}
以 $j = 5$ 為例：$t = 1$，$A[0] = 1$；$i = 4, 3, 2, 1$ 時 $A[i] = 6, 5, 4, 2$ 都比 $1$ 大，依序挪到
$A[5], A[4], A[3], A[2]$；$i = 0$ 時 $1 < A[0] = 1$ 不成立，停，$A[1] = 1$。

條件用的是嚴格的 $<$，和 $t$ 相等的數不會被挪過去，所以這個插入排序是穩定的。 時間和資結書
第~1.3~節一樣，最壞 $\Theta(N^2)$（由大到小排好時），最好 $\Theta(N)$（已經排好時）。
\end{pastsoln}
